\begin{recipebox}{Beefaroni}{8}{30 min}
  \begin{ingredients}
    \item 2 tbsp olive oil
    \item 450\,g beef mince
    \item 1 medium onion, diced
    \item 3 garlic cloves, minced
    \item 1 tbsp onion powder
    \item 2 tsp paprika
    \item \nicefrac{1}{2} tsp chilli powder
    \item 1 tsp cumin
    \item 2 tsp Italian herbs (or dried oregano)
    \item 1 tsp salt
    \item 450\,g dry elbow macaroni
    \item 2 cups tomato sauce (preferably tomato basil)
    \item 2 tbsp tomato paste
    \item 45\,g butter, cubed
    \item 2 cups grated cheese (Cheddar, Jack, Parmesan, or mixed)
    \item Spring onions or parsley, for garnish
  \end{ingredients}

  \begin{method}
    \item Press \textbf{Sauté} and heat the olive oil. Add beef mince and cook for 4–5 minutes until browned, breaking it apart as it cooks.
    \item Add diced onion, garlic, onion powder, paprika, chilli powder, cumin, Italian herbs, and salt; stir well.
    \item Pour in 2 cups of water and scrape the bottom of the pot to deglaze and prevent the \emph{Burn} warning. Cancel \textbf{Sauté}.
    \item Add the pasta on top (do not stir), then spread tomato sauce and tomato paste over it. Scatter butter cubes on top.
    \item Lock the lid and set to \textbf{Manual/Pressure Cook — HIGH} for 4 minutes.
    \item When done, allow natural release for 5 minutes, then perform a quick release.
    \item Open the lid and stir well to combine.
    \item With the pot on \textbf{Keep Warm}, add the cheese and cover with a glass lid for 5 minutes to melt. Alternatively, transfer to a casserole dish and broil to melt cheese.
    \item Garnish with spring onions or parsley and serve hot.
  \end{method}

  \begin{notes}
    \begin{itemize}
      \item Beef mince can be replaced with chicken mince, pork, or lamb.
      \item Any pasta sauce works—tomato basil adds great flavour. If using canned tomatoes instead, add a pinch more salt.
      \item Elbow macaroni is traditional, but penne or fusilli also work well.
      \item Store leftovers in the fridge for 2–3 days. Reheat with a splash of water or milk.
      \item For extra heat, increase the chilli powder or add a pinch of crushed red pepper.
    \end{itemize}
  \end{notes}
\end{recipebox}